\section{Równania stanu}

\subsection{Wektor Stanu}

Jako wektor równań stanu przyjęto:

\begin{equation}
\mathbf{x}=\begin{bmatrix}
    x \\ \dot{x} \\ \theta
\end{bmatrix}
\end{equation}

Ze względu na praktyczną implementację układu sterowania, konieczne było wybranie zmiennych stanu, które mogą być zmierzone fizycznie w rzeczywistym systemie. Wektor stanu zawiera trzy mierzalne wielkości: położenie piłki ($x$), prędkość piłki ($\dot{x}$) oraz kąt belki ($\theta$). To podejście, choć wymaga bardziej zaawansowanego wyprowadzenia matematycznego, gwarantuje fizyczną interpretowalność każdej zmiennej stanu.
Kluczowym pytaniem jest: dlaczego nie stosujemy standardowych metod matematycznych (takich jak postać kanoniczna sterowalna)? Odpowiedź jest prosta — metody kanoniczne generują \textbf{stany matematyczne}, które mogą nie mieć fizycznego odpowiednika. W naszym przypadku chcemy, aby trzecia zmienna stanu ($x_3$) reprezentowała kąt belki ($\theta$), a nie przyspieszenie piłki. 

\subsection{Dekompozycja transmitancji}

Mamy transmitancję całkowitą $G(s)$, która wiąże wyjście (pozycję piłki $x$) z
wejściem (napięciem sterującym $U(s)$). Gdzie, zgodnie z opisem z sekcji
poprzedniej:
\begin{align}
G_{\mathrm{serwo}}(s) &= \frac{\Theta_{\mathrm{serwo}}(s)}{U(s)} = \frac{1}{T_{\mathrm{serwo}} s + 1} \quad \text{(Transmitancja serwa)}\\
G_{\mathrm{ball}}(s) &= \frac{X(s)}{\Theta_{\mathrm{belka}}(s)} = \frac{c_{\mathrm{ball}}}{s^2} \quad \text{(Transmitancja piłki)}
\end{align}

Kluczową rolę odgrywa tu przekładnia mechaniczna $k_{\mathrm{mech}} = \frac{r_{\mathrm{serwo}}}{r_{\mathrm{belka}}} = 0{,}2$, która wiąże kąt serwa $\Theta_{\mathrm{serwo}}(s)$ z kątem belki $\Theta_{\mathrm{belka}}(s)$:

\begin{equation}
\Theta_{\mathrm{belka}}(s) = k_{\mathrm{mech}} \cdot \Theta_{\mathrm{serwo}}(s)
\end{equation}

\subsection{Wyprowadzenie równań różniczkowych dla każdego bloku}

\subsubsection{Blok Serwa}

\begin{equation}
    G_{\mathrm{serwo}}(s) = \frac{\Theta_{\mathrm{serwo}}(s)}{U(s)}=\frac{1}{T_{\mathrm{serwo}} s + 1}
\end{equation}

Mnożymy na krzyż:

\begin{equation}
(T_{\mathrm{serwo}} s + 1) \Theta_{\mathrm{serwo}}(s) = U(s)
\end{equation}

Przechodzimy na dziedzinę czasu:

\begin{equation}
T_{\mathrm{serwo}} \dot{\theta}(t) + \theta(t) = u(t)
\end{equation}

Wyznaczamy najwyższą pochodną ($\dot{\theta}$):

\begin{equation}
\dot{\theta}(t) = \frac{1}{T_{\mathrm{serwo}}} u(t) - \frac{1}{T_{\mathrm{serwo}}} \theta(t)
\label{eq:servo_dynamics}
\end{equation}

\subsubsection{Blok Piłki}

Z rozdziału drugiego wiemy, że dynamika piłki w zależności od kąta belki jest opisana równaniem:

\begin{equation}
\ddot{x}(t) = c_{\mathrm{ball}} \cdot \theta_{\mathrm{belka}}(t)
\label{eq:ball_dynamics}
\end{equation}

gdzie $c_{\mathrm{ball}} = \frac{g}{1 + \frac{J}{m R^2}}$ jest współczynnikiem charakteryzującym dynamikę piłki na belce, a $\theta_{\mathrm{belka}}$ jest kątem belki.

Wiemy, że kąt belki związany jest z kątem serwa przez przekładnię mechaniczną:

\begin{equation}
\theta_{\mathrm{belka}}(t) = k_{\mathrm{mech}} \cdot \theta_{\mathrm{serwo}}(t)
\end{equation}

Stąd:

\begin{equation}
\ddot{x}(t) = c_{\mathrm{ball}} \cdot k_{\mathrm{mech}} \cdot \theta_{\mathrm{serwo}}(t) 
\label{eq:ball_dynamics_with_mech}
\end{equation}

\subsection{Definicja zmiennych stanu}

Chcemy, aby nasz wektor stanu wyglądał tak:

\begin{equation}
x = \begin{bmatrix} x_1 \\ x_2 \\ x_3 \end{bmatrix} = \begin{bmatrix} x \\ \dot{x} \\ \theta \end{bmatrix}
\end{equation}

Gdzie:

\begin{itemize}
\item $x_1 = x$ (Pozycja piłki)
\item $x_2 = \dot{x}$ (Prędkość piłki)
\item $x_3 = \theta$ (Kąt belki)
\end{itemize}

\subsection{Zapis pochodnych zmiennych stanu}

\subsubsection{Pochodna pozycji ($\dot{x}_1$)}

Z definicji prędkości:

\begin{equation}
\dot{x}_1 = \dot{x} = x_2
\end{equation}

\subsubsection{Pochodna prędkości ($\dot{x}_2$)}

To jest przyspieszenie. Bierzemy je z równania \eqref{eq:ball_dynamics_with_mech}, które uwzględnia przekładnię mechaniczną:

\begin{equation}
\ddot{x} = k_{\mathrm{mech}} \cdot c_{\mathrm{ball}} \cdot \theta_{\mathrm{serwo}}
\end{equation}

Ponieważ $\theta_{\mathrm{serwo}} = x_3$ (zmienna stanu), podstawiając:

\begin{equation}
\dot{x}_2 = k_{\mathrm{mech}} \cdot c_{\mathrm{ball}} \cdot x_3
\end{equation}

\subsubsection{Pochodna kąta ($\dot{x}_3$)}

To jest prędkość kątowa serwa. Bierzemy ją z równania \eqref{eq:servo_dynamics}:

\begin{equation}
\dot{x}_3 = \frac{1}{T_{\mathrm{serwo}}} u - \frac{1}{T_{\mathrm{serwo}}} x_3
\end{equation}

\subsection{Postać macierzowa}

Układ równań liniowych:

\begin{equation}
\begin{cases}
\dot{x}_1 = x_2 \\
\dot{x}_2 = k_{\mathrm{mech}} \cdot c_{\mathrm{ball}} \cdot x_3 \\
\dot{x}_3 = \frac{1}{T_{\mathrm{serwo}}} u - \frac{1}{T_{\mathrm{serwo}}} x_3
\end{cases}
\end{equation}

Zapis macierzowy $\dot{x} = Ax + B u$:

\begin{equation}
\begin{bmatrix} \dot{x}_1 \\ \dot{x}_2 \\ \dot{x}_3 \end{bmatrix} = \begin{bmatrix} 0 & 1 & 0 \\ 0 & 0 & k_{\mathrm{mech}} \cdot c_{\mathrm{ball}} \\ 0 & 0 & -\frac{1}{T_{\mathrm{serwo}}} \end{bmatrix} \begin{bmatrix} x_1 \\ x_2 \\ x_3 \end{bmatrix} + \begin{bmatrix} 0 \\ 0 \\ \frac{1}{T_{\mathrm{serwo}}} \end{bmatrix} u
\end{equation}

Oraz równanie wyjścia (mierzymy pozycję $x$):

\begin{equation}
y = \begin{bmatrix} 1 & 0 & 0 \end{bmatrix} \begin{bmatrix} x_1 \\ x_2 \\ x_3 \end{bmatrix}
\end{equation}

